\documentclass[12pt]{article}
\usepackage{amsmath,amsfonts,amssymb}
\usepackage{float}
\usepackage{array}
\usepackage{geometry}
\geometry{
legalpaper, total={177.8mm, 290mm},left=10mm, right=10mm,
top=7mm, bottom=10mm,
}
\setlength\parindent{0pt}
\pagestyle{empty}

\begin{document}

% ---------- Header ----------
\noindent
\begin{minipage}[t]{0.34\textwidth}
\vspace{0pt}
\renewcommand{\arraystretch}{1.3}
\begin{tabular}{|l|p{12mm}|}
\hline
Q/Setter & \\ \hline
H/Dept /Moderator & \\ \hline
\end{tabular}
\end{minipage}%
\begin{minipage}[t]{0.37\textwidth}
\vspace{0pt}
\raggedright
Mymensingh Girls' Cadet College\par
3rd Term End Exam-2026\par
Subject-Statistics First Paper\par
Class- XI\par
\end{minipage}%
\begin{minipage}[t]{0.29\textwidth}
\vspace{0pt}
\raggedleft
\setlength{\fboxsep}{4pt}
\renewcommand{\arraystretch}{1.2}
\fbox{\begin{tabular}{@{}l@{\hspace{6pt}}r@{}}
Subject Code &
\begin{tabular}{|c|c|c|}
\hline 1 & 2 & 9 \\ \hline
\end{tabular}
\end{tabular}}\par
\end{minipage}

\smallskip
\noindent\textbf{Time 2 hours 35 minutes}\hfill\textbf{Marks 50}\par\smallskip

\hrule
\begin{center}
[N.B. -- The figures in the right margin indicate full marks. Read the stems carefully and answer the associated questions. Answer any \textbf{FIVE} questions.]
\end{center}

\begin{enumerate}

% ---------- Chapter 1 ----------
\item
\textbf{The weights (in kg) of eight cadets were recorded in the college hospital. It was later found that the weighing machine had shown 2 kg more than the true weight for every cadet. The recorded weights are:}
\begin{center}
52, 48, 55, 50, 47, 53, 49, 54
\end{center}
\begin{enumerate}
  \item What is change of origin? \hfill 1
  \item Does a change of origin affect the summation of a variable? Explain. \hfill 2
  \item Find $\displaystyle \sum_{i=1}^{8} (x_i-50)$, where $x_i$ is a recorded weight. \hfill 3
  \item Determine the sum of the true weights directly and by using the rules of summation, and compare the results. \hfill 4
\end{enumerate}

% ---------- Chapter 1 ----------
\item
\textbf{A stationery shop sold four types of items in a day. The number of units sold ($x$) and the price per unit in taka ($y$) are given below:}
\begin{center}
\begin{tabular}{l|c|c|c|c}
Item & Pen & Notebook & File & Marker \\ \hline
Units sold ($x$) & 12 & 8 & 5 & 6 \\ \hline
Price per unit ($y$) & 10 & 60 & 30 & 25
\end{tabular}
\end{center}
\begin{enumerate}
  \item What is a discrete variable? \hfill 1
  \item Find $\displaystyle \sum_{i=1}^{4} x_i^2$. \hfill 2
  \item Show, theoretically and with the data, that $\displaystyle \sum_{i=1}^{4} (x_i+y_i) = \sum_{i=1}^{4}x_i + \sum_{i=1}^{4}y_i$. \hfill 3
  \item The shopkeeper claims that the total sales value equals the product of the total units sold and the total of the unit prices. Verify the claim. \hfill 4
\end{enumerate}

% ---------- Chapter 2 ----------
\item
\textbf{A college librarian noted the category of each of the 30 books borrowed on a particular day:}
\begin{center}
History, Fiction, Science, Poetry, Fiction, Poetry,\\
Fiction, Fiction, Science, Fiction, Poetry, History,\\
Biography, Science, Biography, Biography, History, Fiction,\\
Science, Biography, History, Science, Fiction, Science,\\
Fiction, Fiction, History, Science, Fiction, Science
\end{center}
\begin{enumerate}
  \item Which category of book was borrowed the most? \hfill 1
  \item Is the librarian's data primary or secondary? Give a reason. \hfill 2
  \item Draw a bar diagram from the data and explain it. \hfill 3
  \item Would a pie chart represent the data better? Justify by drawing one. \hfill 4
\end{enumerate}

% ---------- Chapter 2 ----------
\item
\textbf{A manufacturer tested 20 LED bulbs and recorded the lifetime of each (in hundreds of hours):}
\begin{center}
42, 55, 61, 48, 53, 67, 58, 50, 63, 45 \\
56, 59, 71, 52, 64, 47, 57, 68, 54, 60
\end{center}
\begin{enumerate}
  \item What is a class interval? \hfill 1
  \item What problem arises if discrete class intervals are used for lifetime data? \hfill 2
  \item Construct a frequency distribution and draw a histogram from it. \hfill 3
  \item Draw an ogive and compare what it shows with what the histogram shows. \hfill 4
\end{enumerate}

% ---------- Chapter 3 ----------
\item
\textbf{Five judges mark a participant in a competition. A more experienced judge's mark carries a larger weight, as shown below:}
\begin{center}
\begin{tabular}{l|c|c|c|c|c}
Judge & 1 & 2 & 3 & 4 & 5 \\ \hline
Mark ($x$) & 2 & 4 & 6 & 7 & 9 \\ \hline
Weight ($w$) & 3 & 2.5 & 2 & 1.5 & 1
\end{tabular}
\end{center}
\begin{enumerate}
  \item Write down the formula of the weighted arithmetic mean. \hfill 1
  \item Why is a weighted mean more suitable than a simple mean here? \hfill 2
  \item Find the weighted arithmetic mean of the marks. \hfill 3
  \item The organisers consider giving all the judges equal weight. Would the participant benefit? Justify with calculations. \hfill 4
\end{enumerate}

% ---------- Chapter 3 ----------
\item
\textbf{The number of bacteria (in thousands) in a culture was counted at the end of five successive hours:}
\begin{center}
\begin{tabular}{l|c|c|c|c|c}
Hour & 1 & 2 & 3 & 4 & 5 \\ \hline
Bacteria ($x$) & 3 & 6 & 12 & 24 & 48
\end{tabular}
\end{center}
\begin{enumerate}
  \item Write down the formula of the geometric mean for ungrouped data. \hfill 1
  \item What happens to the geometric mean if one of the values is 0? Explain. \hfill 2
  \item For $n$ positive values, derive the formula $\displaystyle G = \mbox{antilog}\left(\frac{1}{n}\sum_{i=1}^{n} \log x_i\right)$. \hfill 3
  \item Find the geometric mean of the data directly and by using logarithm and antilogarithm, and compare the results. \hfill 4
\end{enumerate}

\end{enumerate}

\vspace{1cm}

\begin{center}
``You can have data without information, but you cannot have information without data.''\\
-- Daniel Keys Moran
\end{center}

\vfill

\begin{center}
--AAM--
\end{center}

\end{document}
